\section{Logarithmic inversion and fine threshold location}\label{sfh:sec:inverse}
Two distinct inverse problems must be separated. A reversed logarithmic expansion gives an explicit asymptotic size and a practical root-finding seed. The exact-saddle expansion gives a far finer location and, when supplemented with effective bounds or exact adjacent checks, integer thresholds.

\subsection{The exact-saddle smooth function}
For real $x$ sufficiently large let $r_x$ solve $\De C(r_x)=x$, put $b_x=\De^2C(r_x)$, and define
\begin{equation}\label{sfh:eq:F}
 F_J(x)=\log\Gamma(x+1)+C(r_x)-x\log r_x
       -\tfrac12\log(2\pi b_x)+\log S_J(r_x).
\end{equation}
Since $S_J=1+O(r_x/x)$, it is positive eventually. Theorem~\ref{sfh:thm:coeff} gives, at integer $n$,
\begin{equation}\label{sfh:eq:logerror}
 \log A_n=F_J(n)+O_J((r_n/n)^{J+1}).
\end{equation}
Differentiating the saddle relation gives $\dd r_x/\dd x=r_x/b_x$. The implicit derivatives in $C(r_x)-x\log r_x$ cancel, and therefore
\begin{equation}\label{sfh:eq:Fprime}
 F'_J(x)=\psi(x+1)-\log r_x-\frac{\kappa_3(r_x)}{2b_x^2}
       +\frac{r_x}{b_x}\frac{\dd}{\dd r}\log S_J(r_x)
       \sim\log(x/r_x)>0.
\end{equation}
Indeed the explicit rational monomials in \eqref{sfh:eq:Eell} satisfy
$\De E_\ell=O_J(r\eps^\ell)$, and the last contribution is $O_J(\eps/x)$. Thus $F_J$ has a unique inverse on a sufficiently large half-line.

\begin{theorem}[Rounding-safe inverse brackets]\label{sfh:thm:inverse}
Fix $J$. Suppose $K_J$ and a starting point are chosen so that
\[
 |\log A_n-F_J(n)|\le K_J(r_n/n)^{J+1}
\]
for all integers beyond that point. For sufficiently large $M$, define $x_-,x_+$ by
\begin{align}
 F_J(x_-)+K_J(r_{x_-}/x_-)^{J+1}&=\log M,\label{sfh:eq:lowerinverse}\\
 F_J(x_+)-K_J(r_{x_+}/x_+)^{J+1}&=\log M.\label{sfh:eq:upperinverse}
\end{align}
Then the integer threshold $N(M)=\min\{n:A_n\ge M\}$ satisfies
\begin{equation}\label{sfh:eq:ceilings}
 \ceil{x_-}\le N(M)\le\ceil{x_+}.
\end{equation}
If $x_J$ is defined by $F_J(x_J)=\log M$, the bracket width and each endpoint's distance from $x_J$ have order at most
\begin{equation}\label{sfh:eq:inversefine}
 O_J\left(\frac{(r_{x_J}/x_J)^{J+1}}{\log(x_J/r_{x_J})}\right).
\end{equation}
\end{theorem}
\begin{proof}
The two functions on the left of \eqref{sfh:eq:lowerinverse}--\eqref{sfh:eq:upperinverse} are eventually increasing, with derivatives asymptotic to \eqref{sfh:eq:Fprime}. If an integer $n<\ceil{x_-}$, then $n<x_-$, and the upper estimate for $\log A_n$ is below $\log M$. If $n\ge\ceil{x_+}$, its lower estimate is at least $\log M$. The finite early segment is irrelevant for sufficiently large $M$. The mean value theorem and \eqref{sfh:eq:Fprime} give \eqref{sfh:eq:inversefine}. Finally $A_n$ is nondecreasing by adjoining an isolated labelled vertex, and the coefficient formula also implies eventual strict increase.
\end{proof}

The theorem is an asymptotic bracket statement; this article does not supply numerical values for $K_J$ and the starting point. It therefore does not claim that rounding a computed asymptotic inverse alone certifies an exact threshold. When the bracket crosses an integer, exact adjacent coefficient comparisons settle the answer. There is also no need to posit a canonical smooth interpolation of the integer sequence: \eqref{sfh:eq:ceilings} is the precise meaning of fine range location.

\subsection{Three explicit inverse-logarithmic corrections}
Let $y=\log M$, $L=\log y$, $\ell=\log L$, and $c=\log2$.
\begin{theorem}[Reversed logarithmic expansion]\label{sfh:thm:reverse}
For each fixed $J$, the solution $x_J(y)$ of $F_J(x)=y$ satisfies
\begin{align}\label{sfh:eq:reverse}
 x_J(y)=\frac yL\bigg\{1&+\frac{2\ell+1}{L}
 +\frac{4\ell^2-2\ell+c-1}{L^2}\\[-3pt]
 &+\frac{8\ell^3-22\ell^2+8c\ell-c-c^2/2-1/2}{L^3}
 +O\left(\frac{\ell^4}{L^4}\right)\bigg\}.\nonumber
\end{align}
There is a coefficient-by-coefficient recursion for every fixed subsequent logarithmic order.
\end{theorem}
\begin{proof}
Define
\begin{align}
 n_0(r)&=\tfrac12e^rr^2(r+2),\nonumber\\
 g(r)&=r+2\log r-c-1+\log(1+2/r)+\frac1{r+2}.\label{sfh:eq:g}
\end{align}
The exact saddle, Stirling's expansion, and the polynomial perturbation imply
\[
 n=n_0(r)+O(r^4),\qquad F_J(n)=n_0(r)g(r)+O(r^5).
\]
For example, the main factor in the latter expression is
$\log n_0-\log r-1+H(r)/n_0=g(r)$; the saddle perturbation and fixed logarithmic factors fit inside $O(r^5)$. Errors polynomial in $r$ are exponentially small relatively here. They may be discarded for fixed inverse-logarithmic orders, though not for count-relative accuracy. Taking logarithms and then dividing by $g$ yields, to all such orders,
\begin{equation}\label{sfh:eq:logreversal}
 L=r+3\log r-c+\log(1+2/r)+\log g(r),\qquad n=\frac y{g(r)}.
\end{equation}
First,
\[
 g(r)=r+2\log r-c-1+\frac3r-\frac4{r^2}
                  +\frac{20}{3r^3}+O(r^{-4}).
\]
Reversion in \eqref{sfh:eq:logreversal} gives
\begin{align}\label{sfh:eq:rinverse}
 r={}&L-4\ell+c+\frac{14\ell-3c-1}{L}\\
 &+\frac{26\ell^2-12c\ell-54\ell+3c^2/2+12c+7/2}{L^2}
 +O(\ell^3/L^3).\nonumber
\end{align}
Substituting \eqref{sfh:eq:rinverse} in $y/g(r)$ gives \eqref{sfh:eq:reverse}.

For a mechanical arbitrary-order rule set $t=1/L$, regard $\ell,c$ as formal constants, and seek
\[
 R=t^{-1}-4\ell+c+\sum_{k\ge1}a_k(\ell)t^k,
 \qquad \log R=\ell+\log(tR).
\]
Expand \eqref{sfh:eq:logreversal} with $\log g(R)=\ell+\log(tg(R))$. At power $t^k$ the new unknown $a_k$ occurs with coefficient one in $R$, while its appearances in logarithms are delayed by at least one power. Solve successively, then expand $1/[t g(R)]$. This gives polynomials $q_k(\ell)$ of degree at most $k$ in
\[
 n=\frac yL\sum_{k\ge0}\frac{q_k(\ell)}{L^k}.
\]
Taylor's theorem for $\log(1+u)$, where $|u|=O(\ell/L)$, bounds a truncation through relative order $m$ by $O_m(\ell^{m+1}/L^{m+1})$. The derivative of $\log(n_0(r)g(r))$ is $1+O(1/r)$, so inversion preserves that order. This justifies the formal recursion as an asymptotic expansion.
\end{proof}

The $\log2$ terms record the factor $1/2$ in $H$. Even the error in \eqref{sfh:eq:reverse} is eventually much larger than one. A leading Lambert-$W$ seed $x\approx y/W(y)$ is still coarser, and neither seed should be substituted into a counting formula and called relatively accurate.

\subsection{Newton refinement to the exact-saddle scale}
Use \eqref{sfh:eq:reverse} to initialize Newton iteration for the full function $F_J$, with derivative \eqref{sfh:eq:Fprime}. On $|x-x_J|\le x_J/2$, direct differentiation gives $F'_J\asymp L$ and $F''_J=O(1/x_J)$. Taylor's theorem therefore yields
\[
 |e_{k+1}|\le\frac{C|e_k|^2}{x_JL}.
\]
The logarithmic seed has vanishing relative error and eventually lies in this basin. Once the numerical residual is at most the coefficient-error scale $(r/x)^{J+1}$, the mean value theorem converts it to \eqref{sfh:eq:inversefine}, subject to the numerical error margin. Solving the reduced equation $n_0g=y$ alone would not obtain this finer accuracy. Effective integer certification still requires the additional bounds or exact checks specified above.

\begin{remark}[{[write]} The inverse and the transseries volume, 7~October 2026]\label{sfh:rem:transseries}
Compared statement by statement with the repository's transseries
volume~\cite{TSvol}; the volume's symbols are not used here.
\begin{enumerate}\raggedright
\item \emph{The growth.} $\log A_n=n\,g(r)+O(r^5)$ with $g(r)\sim r\sim\log n$
contains the term $n\log n$, so the sequences are not of the
exponential--power form of \file{p0:def:model}; no instance of the volume's
model is claimed.
\item \emph{The reversed logarithmic expansion.} Writing
$\log g(r)=\log r+\log(g(r)/r)$ with
$g(r)/r=1+(2\log r-c-1)/r+3/r^2-\cdots$, the reduced
equation~\eqref{sfh:eq:logreversal} for $r$ is literally the volume's
monomial--logarithmic equation (\file{plt:def:lw-monomial-log-datum}) in the
unknown $r$: unit slope, logarithmic coefficient~$4$, constant $-\log2$, and a
perturbation in $1/r$ of positive order with coefficients polynomial in
$\log r$. The three hypotheses of \file{plt:thm:lw-template} hold, so the
expansion~\eqref{sfh:eq:rinverse} of $r$ (re-expanded from the template's
variable $L+\log2$ into powers of $1/L$) is an exact formal instance of the
template, and~\eqref{sfh:eq:reverse} follows from it by the explicit
substitution $x=y/g(r)$. The write's own reversion reproduces both. The
reduction itself, valid to every inverse-logarithmic order, is the source's
argument in the proof of Theorem~\ref{sfh:thm:reverse}.
\item \emph{The exact-saddle inverse.} $F_J$ differs from the reduced model
by terms polynomial in $r$, relatively exponentially small in $r$: below every
order of the template's chart. So $x_J$ beyond~\eqref{sfh:eq:reverse}, the
Newton refinement and the brackets of Theorem~\ref{sfh:thm:inverse} are not
instances of \file{plt:thm:lw-template}; they are proved directly, by the
mean value theorem for the explicit $F_J$.
\item \emph{The thresholds.} $N(M)$ is the volume's staircase
(\file{p0:def:three-inverses}(1)) of the eventually increasing $A_n$. The
functions $F_J\pm K_J(r_x/x)^{J+1}$ do not take the values $\log A_n$ at the
integers, so they are not admissible interpolations, and
Theorem~\ref{sfh:thm:inverse}, proved directly at the integers, is an
analogue of \file{p0:thm:staircase}(2), not an instance; the source's refusal
to posit a canonical smooth interpolation is the distinction of
\file{p0:def:three-inverses}.
\end{enumerate}
\end{remark}
